\documentclass[aspectratio=169]{beamer}

% ==== Theme & basic look ====
\usetheme{Madrid}
\usecolortheme{seahorse}
\setbeamertemplate{navigation symbols}{}
\setbeamertemplate{footline}[frame number]


% Packages for mathematics  
\usepackage{amsmath}  
\usepackage{amsfonts}  
\usepackage{amssymb}  
\usepackage{CJKutf8}
\usepackage{algorithm}
\usepackage{algorithmic}
% Support for \Tilde macro used throughout
\providecommand{\Tilde}[1]{\tilde{#1}}
\providecommand{\proofname}{Proof}
\newenvironment<>{proofs}[1][\proofname]{%
    \par
    \def\insertproofname{#1.}%
    \usebeamertemplate{proof begin}#2}
  {\usebeamertemplate{proof end}}
\makeatother

% Math shorthand and vector symbols
\newcommand{\grad}{\nabla}
\newcommand{\EE}{\mathbb{E}}
\newcommand{\dd}{\mathrm{d}}
\newcommand{\kbT}{k_B T}
\newcommand{\vx}{\mathbf{x}}
\newcommand{\eps}{\epsilon}
\newcommand{\z}{\mathbf{z}}
\newcommand{\vinst}{\mathbf{v}}
\newcommand{\uavg}{\overline{\mathbf{u}}}
\newcommand{\dv}[2]{\frac{d#1}{d#2}}
\newcommand{\td}{\frac{d}{dt}}
\newcommand{\x}{\mathbf{x}}
\newcommand{\coloneqq}{\mathrel{\vcenter{\baselineskip0.5ex\lineskiplimit0pt\hbox{\scriptsize.}\hbox{\scriptsize.}}}=}
\newcommand{\sg}{\mathrm{sg}}
\newcommand{\JVP}{\mathrm{JVP}}

% Bold vector/time-indexed variables
\newcommand{\vxz}{\mathbf{x}_0}
\newcommand{\vxt}{\mathbf{x}_t}
\newcommand{\alpt}{\alpha_t}
\newcommand{\sigt}{\sigma_t}
\newcommand{\veps}{\boldsymbol{\epsilon}}
\newcommand{\vF}{\mathbf{F}}

% Title page details:   
\title[Understanding Flow Matching]{Understanding Flow Matching from a mathematical view}  
\author{Pipi Hu}  
\institute[BIMSA]{Beijing Institute of Mathematical Sciences and Applications}  
\date{\today}  

\pgfdeclareimage[width=2cm]{logo}{../figures/logo.png}
\logo{\pgfuseimage{logo}}

% Outline slide (moved to later; avoid frame before \begin{document})

\setbeamertemplate{navigation symbols}{}
\AtBeginSection[]  
{  
    \begin{frame}<beamer>{Outline}         
        \tableofcontents[currentsection]  
    \end{frame}  
}  
  
\AtBeginSubsection[]  
{  
    \begin{frame}<beamer>{Outline}         
        \tableofcontents[currentsection,currentsubsection]  
    \end{frame}  
}  

\begin{document}  
  
% Title slide  
\begin{frame}  
  \titlepage  
\end{frame}  

\begin{frame}{Feynman's words}
    \centering
        What I cannot create, I do not understand.
    \end{frame}
    
\begin{frame}{Mathematical motivation for Flow Matching}
\textbf{The fundamental challenge}: Learn a velocity field \(\mathbf{v}(\vx,t)\) that transforms a simple base distribution \(p_1(\vx)\) into the data distribution \(p_0(\vx)\) through the continuity equation:
\begin{equation}
    \partial_t p(\vx,t) + \grad\cdot\big(p(\vx,t)\,\mathbf{v}(\vx,t)\big) = 0.
\end{equation}

\textbf{Key advantages over score-based methods}:
\begin{itemize}
    \item \textbf{No score estimation}: Avoids computing \(\grad_{\vx}\log p(\vx,t)\) during training
    \item \textbf{Analytic supervision}: Provides closed-form velocity targets via conditional paths
    \item \textbf{Deterministic sampling}: Enables ODE integration with large solver steps
    \item \textbf{Mathematical elegance}: Direct connection to optimal transport theory
\end{itemize}
\end{frame}
  
% Presentation slides  

\section{Continuity equation and probability flow}

\begin{frame}{The continuity equation: mathematical foundation}
\textbf{Theorem}: Let \(\mathbf{v}(\vx,t)\) be a smooth velocity field. Then the density \(p(\vx,t)\) transported by the ODE \(\dot{\vx} = \mathbf{v}(\vx,t)\) satisfies the continuity equation:
\begin{equation}
    \partial_t p(\vx,t) + \grad\cdot\big(p(\vx,t)\,\mathbf{v}(\vx,t)\big) = 0.
\end{equation}

\textbf{Proof sketch}: For any test function \(\phi(\vx)\), we have:
\begin{align}
    \frac{d}{dt}\int \phi(\vx) p(\vx,t) d\vx &= \int \phi(\vx) \partial_t p(\vx,t) d\vx \\
    &= \int \phi(\vx) \mathbf{v}(\vx,t) \cdot \grad p(\vx,t) d\vx \\
    &= -\int \grad\phi(\vx) \cdot \mathbf{v}(\vx,t) p(\vx,t) d\vx
\end{align}
The last equality follows from integration by parts and the chain rule.
\end{frame}

\begin{frame}{Derivation: target velocity for Gaussian path}
Given \(\vxt=\boldsymbol{\mu}_t(\vxz)+\sigt\veps\) with \(\veps\sim\mathcal{N}(0,\mathbf{I})\),
\begin{equation}
    \frac{\dd}{\dd t}\,\vxt = \dot{\boldsymbol{\mu}}_t(\vxz) + \dot{\sigt}\,\veps.
\end{equation}
Condition on \(\vxt=\vx\) and \(\vxz\). Since \(\veps=(\vx-\boldsymbol{\mu}_t(\vxz))/\sigt\),
\begin{align}
    \mathbf{u}_t(\vx\mid\vxz)
    &= \EE\big[\tfrac{\dd}{\dd t}\,\vxt\mid \vxt=\vx,\,\vxz\big]\\
    &= \dot{\boldsymbol{\mu}}_t(\vxz) + \dot{\sigt}\, \EE[\veps\mid\vxt=\vx,\,\vxz]\\
    &= \dot{\boldsymbol{\mu}}_t(\vxz) + \dot{\sigt}\,\frac{\vx-\boldsymbol{\mu}_t(\vxz)}{\sigt}\\
    &= \dot{\boldsymbol{\mu}}_t(\vxz) + \frac{\dot{\sigt}}{\sigt}\Big(\vx-\boldsymbol{\mu}_t(\vxz)\Big).
\end{align}
\end{frame}

\begin{frame}{Conditional probability paths: mathematical framework}
\textbf{Definition}: Given a data distribution \(p_{\text{data}}(\vxz)\) and base distribution \(p_{\text{base}}(\vx)\), we define conditional paths \(q_t(\vx\mid\vxz)\) such that:
\begin{equation}
    p_t(\vx) = \int p_{\text{data}}(\vxz)\; q_t(\vx\mid\vxz)\, \dd\vxz, \quad t\in[0,1],
\end{equation}
with boundary conditions \(p_0=p_{\text{data}}\) and \(p_1=p_{\text{base}}\).

\textbf{Key insight}: The conditional velocity field
\begin{equation}
    \mathbf{u}_t(\vx\mid\vxz) := \EE\big[\tfrac{\dd}{\dd t}\, \vxt \;\big|\; \vxt=\vx,\,\vxz\big]
\end{equation}
provides \emph{analytic supervision} for learning the marginal velocity field \(\mathbf{v}(\vx,t)\).

\textbf{Mathematical advantage}: No need to estimate \(\grad_{\vx}\log p(\vx,t)\) during training!
\end{frame}

\section{Gaussian conditional paths and targets}

\begin{frame}{Gaussian conditional paths: mathematical derivation}
\textbf{Assumption}: Gaussian conditional paths
\begin{equation}
    q_t(\vx\mid\vxz) = \mathcal{N}\big(\vx;\; \boldsymbol{\mu}_t(\vxz),\, \sigma_t^2\mathbf{I}\big),\quad \vxt = \boldsymbol{\mu}_t(\vxz) + \sigt\veps.
\end{equation}

\textbf{Mathematical derivation}: Taking the time derivative of the path:
\begin{equation}
    \frac{\dd}{\dd t}\,\vxt = \dot{\boldsymbol{\mu}}_t(\vxz) + \dot{\sigt}\,\veps,\quad \veps=\frac{\vxt-\boldsymbol{\mu}_t(\vxz)}{\sigt}.
\end{equation}

\textbf{Conditional expectation}: Since \(\veps\) is independent of \(\vxt\) given \(\vxz\), we have:
\begin{align}
    \mathbf{u}_t(\vx\mid\vxz) &= \EE\big[\tfrac{\dd}{\dd t}\,\vxt \;\big|\; \vxt=\vx,\,\vxz\big] \\
    &= \dot{\boldsymbol{\mu}}_t(\vxz) + \dot{\sigt}\,\EE[\veps\mid\vxt=\vx,\,\vxz] \\
    &= \dot{\boldsymbol{\mu}}_t(\vxz) + \dot{\sigt}\,\frac{\vx-\boldsymbol{\mu}_t(\vxz)}{\sigt}
\end{align}

\textbf{Final result}:
\begin{equation}
    \boxed{\;\mathbf{u}_t(\vx\mid\vxz) = \dot{\boldsymbol{\mu}}_t(\vxz) + \frac{\dot{\sigt}}{\sigt}\Big(\vx-\boldsymbol{\mu}_t(\vxz)\Big).\;}
\end{equation}
\end{frame}

\begin{frame}{Mathematical analysis of path choices}
\textbf{1. Linear bridge} (optimal transport):
\begin{itemize}
    \item \(\boldsymbol{\mu}_t(\vxz)=(1-t)\,\vxz\), \(\sigt=t\,\sigma_{\max}\)
    \item \textbf{Mathematical property}: Minimizes Wasserstein-2 distance
    \item \textbf{Velocity}: \(\mathbf{u}_t(\vx\mid\vxz) = \frac{\sigma_{\max}}{t}(\vx-\vxz)\)
\end{itemize}

\textbf{2. VP-like} (DDPM-inspired):
\begin{itemize}
    \item \(\boldsymbol{\mu}_t(\vxz)=\alpha(t)\,\vxz\), \(\sigt=\sqrt{1-\alpha^2(t)}\)
    \item \textbf{Mathematical property}: Preserves energy structure
    \item \textbf{Velocity}: \(\mathbf{u}_t(\vx\mid\vxz) = \dot{\alpha}(t)\vxz + \frac{\dot{\sigt}}{\sigt}(\vx-\alpha(t)\vxz)\)
\end{itemize}

\textbf{3. VE-like} (SMLD-inspired):
\begin{itemize}
    \item \(\boldsymbol{\mu}_t(\vxz)=\vxz\), \(\sigt=\sigma(t)\) increasing
    \item \textbf{Mathematical property}: Preserves data structure
    \item \textbf{Velocity}: \(\mathbf{u}_t(\vx\mid\vxz) = \frac{\dot{\sigma}(t)}{\sigma(t)}(\vx-\vxz)\)
\end{itemize}

\textbf{Key insight}: All choices provide \emph{closed-form} \(\mathbf{u}_t\) via the boxed formula!
\end{frame}

\begin{frame}{Derivation: VE/VP target forms}
\small
\textbf{VE-like.} \(\boldsymbol{\mu}_t(\vxz)=\vxz\), \(\sigt=\sigma(t)\). Then
\begin{equation}
    \mathbf{u}_t = 0 + \frac{\dot\sigt}{\sigt}\,(\vx-\vxz).
\end{equation}
\textbf{VP-like.} \(\boldsymbol{\mu}_t(\vxz)=\alpha(t)\,\vxz\), \(\sigt=\sqrt{1-\alpha^2(t)}\). Using
\(\mathbf{u}_t=\dot{\boldsymbol{\mu}}_t + (\dot\sigt/\sigt)(\vx-\boldsymbol{\mu}_t)\), we get
\begin{align}
    \mathbf{u}_t(\vx\mid\vxz)
    &= \dot\alpha\,\vxz + \frac{\dot\sigt}{\sigt}\,\vx - \frac{\dot\sigt}{\sigt}\,\alpha\,\vxz\\
    &= \frac{\dot\sigt}{\sigt}\,\vx + \Big(\dot\alpha - \alpha\,\frac{\dot\sigt}{\sigt}\Big)\,\vxz.
\end{align}
No explicit form for \(\dot\sigt/\sigt\) is required to train; schedules determine it numerically.
\end{frame}

\section{Training and inference}

\begin{frame}{Flow Matching: mathematical training objective}
\textbf{Theoretical foundation}: Given conditional paths \(q_t(\vx\mid\vxz)\), the optimal velocity field satisfies:
\begin{equation}
    \mathbf{v}^*(\vx,t) = \EE_{\vxz\mid\vxt=\vx}\big[\mathbf{u}_t(\vx\mid\vxz)\big]
\end{equation}

\textbf{Training procedure}:
\begin{enumerate}
    \item Sample \(t\sim\mathcal{U}[0,1]\), \(\vxz\sim p_{\text{data}}\), \(\vxt\sim q_t(\cdot\mid\vxz)\)
    \item Compute analytic target \(\mathbf{u}_t(\vxt\mid\vxz)\) from chosen path
    \item Minimize regression loss
\end{enumerate}

\textbf{Mathematical objective}:
\begin{equation}
    J_{\text{FM}} = \EE\big[\,\|\, \mathbf{v}_\theta(\vxt,t) - \mathbf{u}_t(\vxt\mid\vxz)\,\|^2\,\big].
\end{equation}

\textbf{Key mathematical advantage}: No score estimation needed; all targets are analytic!
\end{frame}

\begin{frame}{Mathematical optimality of Flow Matching}
\textbf{Theorem}: The population minimizer of the Flow Matching objective
\begin{equation}
    J_{\text{FM}}(\theta)=\EE\big[\|\mathbf{v}_\theta(\vxt,t)-\mathbf{u}_t(\vxt\mid\vxz)\|^2\big]
\end{equation}
is the posterior-averaged velocity field:
\begin{equation}
    \mathbf{v}^*(\vx,t)=\EE\big[\,\mathbf{u}_t(\vxt\mid\vxz)\,\big|\,\vxt=\vx,\,t\big].
\end{equation}

\textbf{Mathematical proof}:
\begin{proofs}
Let \(Y=\mathbf{u}_t(\vxt\mid\vxz)\), \(X=(\vxt,t)\). For any predictor \(g(X)\),
\begin{align}
\EE[\|Y-g(X)\|^2]&=\EE[\|Y-\EE[Y\mid X]\|^2]+\EE[\|g(X)-\EE[Y\mid X]\|^2] \\
&\quad + 2\EE[(Y-\EE[Y\mid X])\cdot(g(X)-\EE[Y\mid X])]
\end{align}
The cross term vanishes by the tower property, and the second term is minimized when \(g(X)=\EE[Y\mid X]\).
\end{proofs}

\textbf{Key insight}: Flow Matching learns the \emph{optimal} velocity field that respects the continuity equation!
\end{frame}

\begin{frame}{Mathematical foundation of inference}
\textbf{Theoretical basis}: The learned velocity field \(\mathbf{v}_\theta(\vx,t)\) defines a deterministic ODE:
\begin{equation}
    \frac{\dd}{\dd t}\,\vx = \mathbf{v}_\theta(\vx,t), \quad t\in[0,1]
\end{equation}

\textbf{Mathematical properties}:
\begin{itemize}
    \item \textbf{Existence}: Under Lipschitz conditions, solutions exist and are unique
    \item \textbf{Consistency}: The ODE preserves the continuity equation
    \item \textbf{Stability}: Deterministic integration avoids accumulation of stochastic errors
\end{itemize}

\textbf{Generation process}:
\begin{enumerate}
    \item Initialize \(\vx(1)\sim p_1\) (e.g., \(\mathcal{N}(0,\mathbf{I})\))
    \item Integrate ODE backwards: \(\frac{\dd}{\dd t}\,\vx = \mathbf{v}_\theta(\vx,t)\), \(t:1\to 0\)
    \item Use standard ODE solvers (Euler, RK4, Dormand-Prince, etc.)
\end{enumerate}

\textbf{Mathematical advantages}:
\begin{itemize}
    \item Deterministic sampling (reproducible)
    \item Fast with large solver steps
    \item No stochastic correctors needed
\end{itemize}
\end{frame}

\section{Relation to diffusion and score models}

\begin{frame}{Mathematical connection to diffusion models}
\textbf{Theoretical link}: For a forward SDE \(\dd\vx = f(\vx,t)\,\dd t + g(t)\,\dd B_t\), the probability flow ODE is:
\begin{equation}
    \frac{\dd}{\dd t}\,\vx = f(\vx,t) - \tfrac{1}{2}g^2(t)\,\grad_{\vx}\log p(\vx,t).
\end{equation}

\textbf{Key mathematical insight}: When Flow Matching uses Gaussian paths that match the SDE's marginal distributions \(p_t\), the learned velocity field \(\mathbf{v}_\theta\) approximates the probability flow drift:
\begin{equation}
    \mathbf{v}_\theta(\vx,t) \approx f(\vx,t) - \tfrac{1}{2}g^2(t)\,\grad_{\vx}\log p(\vx,t)
\end{equation}

\textbf{Crucial advantage}: Flow Matching achieves this approximation \emph{without} ever computing the score \(\grad_{\vx}\log p(\vx,t)\) during training!
\end{frame}

\begin{frame}{Mathematical derivation: probability flow ODE}
\textbf{Setup}: Forward SDE \(\dd\vx=f(\vx,t)\,\dd t + g(t)\,\dd B_t\) with density \(p(\vx,t)\) solving the Fokker-Planck equation:
\begin{equation}
    \partial_t p = -\grad\cdot(f p) + \tfrac{1}{2}g^2\,\triangle p.
\end{equation}

\textbf{Key insight}: A deterministic ODE \(\dot{\vx}=\mathbf{v}(\vx,t)\) transports densities via the continuity equation:
\begin{equation}
    \partial_t p + \grad\cdot(p\,\mathbf{v}) = 0.
\end{equation}

\textbf{Mathematical derivation}: Equating the two equations:
\begin{align}
    -\grad\cdot(p\,\mathbf{v}) &= -\grad\cdot(f p) + \tfrac{1}{2}g^2\,\triangle p\\
    \Rightarrow\; \mathbf{v} &= f - \tfrac{1}{2}g^2\,\frac{\grad p}{p} 
    \;=\; f - \tfrac{1}{2}g^2\,\grad\log p.
\end{align}

\textbf{Result}: The probability flow ODE drift is \(\mathbf{v} = f - \tfrac{1}{2}g^2\,\grad\log p\).
\end{frame}

\begin{frame}{Specializing VE/VP-style paths}
\small
\begin{itemize}
    \item VE-style: \(\boldsymbol{\mu}_t(\vxz)=\vxz\). Target becomes \(\mathbf{u}_t=\tfrac{\dot\sigt}{\sigt}(\vx-\vxz)\).
    \item VP-style: \(\boldsymbol{\mu}_t=\alpha(t)\vxz\), \(\sigt=\sqrt{1-\alpha^2(t)}\). Target is \(\mathbf{u}_t=\dot\alpha\,\vxz + \tfrac{\dot\sigt}{\sigt}(\vx-\alpha\vxz)\).
\end{itemize}
Averaging over \(\vxz\) under the path recovers the probability flow drift forms used by diffusion models.
\end{frame}

\section{Practical recipe}

\begin{frame}{Complete training recipe}
\textbf{Step-by-step procedure}:
\begin{enumerate}
    \item \textbf{Sample data}: \(\vxz\sim p_{\text{data}}\), \(t\sim\mathcal{U}[0,1]\), \(\veps\sim \mathcal{N}(0,\mathbf{I})\)
    \item \textbf{Forward pass}: \(\vxt = \boldsymbol{\mu}_t(\vxz) + \sigt\veps\)
    \item \textbf{Compute target}: \(\mathbf{u}_t(\vxt\mid\vxz) = \dot{\boldsymbol{\mu}}_t(\vxz) + (\dot\sigt/\sigt)(\vxt-\boldsymbol{\mu}_t(\vxz))\)
    \item \textbf{Regression loss}: \(\|\mathbf{v}_\theta(\vxt,t)-\mathbf{u}_t(\vxt\mid\vxz)\|^2\)
\end{enumerate}

\textbf{Implementation notes}:
\begin{itemize}
    \item Use standard optimizers (Adam, AdamW)
    \item Typical learning rates: \(10^{-4}\) to \(10^{-3}\)
    \item Training time similar to diffusion models
\end{itemize}
\end{frame}

\begin{frame}{Complete inference recipe}
\textbf{Generation procedure}:
\begin{enumerate}
    \item \textbf{Initialize}: \(\vx(1)\sim p_1\) (e.g., \(\mathcal{N}(0,\mathbf{I})\))
    \item \textbf{ODE integration}: \(\dot{\vx}=\mathbf{v}_\theta(\vx,t)\) from \(t=1\) to \(t=0\)
    \item \textbf{Optional}: Add predictor-corrector steps for higher quality
\end{enumerate}

\textbf{Solver recommendations}:
\begin{itemize}
    \item \textbf{Fast}: Euler method (1-10 steps)
    \item \textbf{Accurate}: RK4, Dormand-Prince (10-50 steps)
    \item \textbf{Adaptive}: Use error control for automatic step sizing
\end{itemize}
\end{frame}

\section{MeanFlow: Average Velocity Field}
\begin{frame}{Motivation}
    \begin{itemize}
    \item Goal: map a simple prior (e.g., $\eps\sim \mathcal{N}(0,I)$) to the data distribution $p_{\text{data}}$.
    \item Diffusion / Flow Matching usually require \emph{multi-step} ODE/SDE integration during sampling.
    \item Desire: \textbf{single-step (1-NFE)} generation while retaining high fidelity and stability.
    \item Key idea in MeanFlow: learn a \emph{ground-truth average velocity field} that summarizes an entire time interval.
    \end{itemize}
    \end{frame}

\begin{frame}{Rectified / Flow Matching setup}
Given schedules $a_t, b_t$ and time $t\in[0,1]$:
\begin{align}
\z_t &= a_t\,\x + b_t\,\eps, \\
\vinst_t &\equiv \dv{\z_t}{t} = a'_t\,\x + b'_t\,\eps. \label{eq:cond-vel}
\end{align}
\vspace{0.5em}
\textbf{Marginal velocity} (ground-truth field for FM):
\begin{equation}
\vinst(\z_t, t) \;=\; \EE_{p_t(\vinst_t\mid \z_t)}\big[\vinst_t\big].
\end{equation}
\textbf{Training (conditional FM):}
\begin{equation}
\mathcal{L}_{\mathrm{CFM}} = \EE_{t,\x,\eps}\,\big\| v_\theta(\z_t,t) - \vinst_t\big\|_2^2, \quad \text{with } \z_t = a_t\x+b_t\eps.
\end{equation}
\end{frame}


\begin{frame}{Sampling in Flow Matching}
The probability-flow ODE:
\begin{equation}
\dv{\z_t}{t} = \vinst(\z_t,t), \qquad \text{(start from }\z_1=\eps\text{)}.
\end{equation}
With a discrete solver (Euler): $\;\z_{t_{i+1}}\approx \z_{t_i} + (t_{i+1}-t_i)\,\vinst(\z_{t_i},t_i)$.
\begin{block}{Issue}
Coarse discretization along \emph{curved} trajectories leads to significant numerical error $\Rightarrow$ many steps are needed.
\end{block}
\end{frame}


\begin{frame}{Definition: Average Velocity}
For any pair $r<t$ and point $\z_t$ along the flow path:
\begin{equation}\label{eq:u-def}
\uavg(\z_t, r, t) \;\coloneqq\; \frac{1}{t-r} \int_{r}^{t} \vinst(\z_\tau,\tau)\, \dd\tau.
\end{equation}
\begin{itemize}
\item $\uavg$ is an \textbf{underlying field} induced by $\vinst$ (no networks yet).
\item As $r\to t$, \; $\uavg(\z_t,r,t) \to \vinst(\z_t,t)$.
\item Composition over intervals holds: one big step equals two smaller consecutive steps (additivity of integrals).
\end{itemize}
\end{frame}

\begin{frame}{MeanFlow Identity (core relation)}
Starting from Eq.~\eqref{eq:u-def}:
\begin{equation}
(t-r)\,\uavg(\z_t,r,t) \,=\, \int_r^t \vinst(\z_\tau,\tau)\,\dd\tau.
\end{equation}
Differentiate w.r.t. $t$ (treating $r$ as constant) and use the fundamental theorem of calculus:
\begin{equation}\label{eq:meanflow-identity}
\boxed{\;\uavg(\z_t,r,t) \,=\, \vinst(\z_t,t) \, - \, (t-r)\,\td\,\uavg(\z_t,r,t)\;}
\end{equation}
where the total derivative expands to
\begin{equation}\label{eq:total-deriv}
\td\,\uavg(\z_t,r,t) \;=\; \underbrace{\vinst(\z_t,t)}_{\dv{\z_t}{t}}\!\cdot \partial_{\z}\uavg(\z_t,r,t) \; +\; \partial_t\uavg(\z_t,r,t).
\end{equation}
\vspace{0.3em}
\textit{Interpretation:} Eq.~\eqref{eq:meanflow-identity} links instantaneous and average velocities \emph{without} introducing networks.
\end{frame}


\begin{frame}{Trainable objective}
Parameterize a network $\uavg_\theta(\z, r, t)$. Use the MeanFlow identity to form a target (replace $\vinst$ by sample-conditional $\vinst_t$):
\begin{align}
\text{Target:}\quad \uavg_{\text{tgt}} \;&=\; \vinst_t \; -\; (t-r)\,\big(\underbrace{\vinst_t \cdot \partial_{\z}\uavg_\theta}_{\text{JVP}} + \partial_t\uavg_\theta\big). \\
\mathcal{L}(\theta) \;&=\; \EE\,\big\|\,\uavg_\theta(\z_t, r, t) - \sg(\uavg_{\text{tgt}})\,\big\|_2^2.
\end{align}
\begin{itemize}
\item $\sg(\cdot)$: stop-gradient; avoids higher-order differentiation.
\item JVP (Jacobian-vector product) computed efficiently via autodiff.
\item If $r=t$, the correction term vanishes $\Rightarrow$ reduces to standard Flow Matching.
\end{itemize}
\end{frame}


\begin{frame}{Pseudocode (training)}
\textbf{MeanFlow training algorithm}:
\begin{enumerate}
\item Sample $t,r\sim \text{LogitNormal}$; set $t\leftarrow\max(r,t)$, $r\leftarrow\min(r,t)$
\item Sample $\x\sim p_{\text{data}}$, $\eps\sim\mathcal{N}(0,I)$; set $\z_t=(1-t)\x + t\eps$, $\vinst_t=\eps-\x$
\item Compute $\big(\uavg_\theta,\; \td\uavg_\theta\big)=\JVP\big(\uavg_\theta(\cdot,r,t), \; (\vinst_t,0,1)\big)$
\item Set $\uavg_{\text{tgt}}\leftarrow \vinst_t - (t-r)\,\td\uavg_\theta$
\item Minimize $\;\|\uavg_\theta(\z_t,r,t)-\sg(\uavg_{\text{tgt}})\|_2^2$ (optionally with adaptive weights)
\end{enumerate}
\end{frame}


\begin{frame}{Sampling (1 step or few steps)}
Replace the time integral by the learned average velocity:
\begin{equation}
\z_r \;=\; \z_t - (t-r)\,\uavg_\theta(\z_t,r,t).
\end{equation}
\textbf{1-NFE sampling}: take $(r,t)=(0,1)$, $\z_1=\eps$,
\begin{equation}
\boxed{\;\x \equiv \z_0 = \eps - \uavg_\theta(\eps,0,1)\;}\;.
\end{equation}
\begin{itemize}
\item Few-step variants are immediate by chaining intervals.
\end{itemize}
\end{frame}

\begin{frame}{CFG inside the field (still 1-NFE)}
Define a CFG velocity field
\begin{equation}
\vinst_{\text{cfg}}(\z_t,t\mid c) \;=\; \omega\,\vinst(\z_t,t\mid c) + (1-\omega)\,\vinst(\z_t,t).
\end{equation}
\textbf{Key insight}: CFG can be applied directly to the average velocity field:
\begin{equation}
\uavg_{\text{cfg}}(\z_t,r,t\mid c) \;=\; \omega\,\uavg(\z_t,r,t\mid c) + (1-\omega)\,\uavg(\z_t,r,t).
\end{equation}
This maintains the 1-NFE property while enabling conditional generation.
\end{frame}

\section{References}

\begin{frame}{Key references and resources}
\textbf{Core papers}:
\begin{itemize}
    \item Lipman et al., \textit{Flow Matching: A unifying perspective of score-based training and generative modeling} (Guide and Code, 2024)
    \item Song et al., \textit{Score-based generative modeling through stochastic differential equations} (2021)
    \item Ho et al., \textit{Denoising Diffusion Probabilistic Models} (2020)
\end{itemize}

\textbf{Implementation resources}:
\begin{itemize}
    \item Flow Matching Guide and Code (GitHub)
    \item Open source libraries: diffusers, torchdyn
    \item Tutorial notebooks and examples
\end{itemize}
\end{frame}

% Thank you slide  
\begin{frame}{Welcome inputs}  
  \centering  
  Any comments?
  \\[2em]
  Welcome your inputs!  
\end{frame}  
  
\end{document}  


